\documentclass[../../../include/open-logic-section]{subfiles}

\begin{document}

\olfileid{sth}{ordinals}{iso}
\olsection{تماثلات الترتيب}

نبدأ بطريقة نقارن بها بين الترتيبات الجيدة؛ فبها يتبين \emph{مدى}
ثبات بنيتها.

\begin{defn}
\emph{الترتيب الجيد} هو الزوج~$\tuple{A, <}$ إذا كانت~$<$ ترتب
$A$ ترتيبًا جيدًا. ونقول إن الترتيبين الجيدين~$\tuple{A, <}$
و$\tuple{B, \lessdot}$ \emph{متماثلان ترتيبيًا} إذا وفقط إذا وجد
!!a{bijection}~$f \colon A \to B$ يحقق الشرط الآتي: $x < y$ إذا
وفقط إذا كان~$f(x) \lessdot f(y)$. وحينئذ نكتب
$\ordeq{\tuple{A, <}}{\tuple{B, \lessdot}}$، ونسمي~$f$
\emph{تماثلًا ترتيبيًا}.
\end{defn}
\noindent
ونقتصر فيما يلي على لفظ «التماثل» اختصارًا لـ«تماثل الترتيب». والمعنى
الحدسي أن التماثلات !!{bijection}s تحفظ البنية. ولنذكر أولًا بعض
خواصها اليسيرة.

\begin{lem}\ollabel{isoscompose}
إذا رُكّبت التماثلات كان تركيبها تماثلًا؛ فإذا كان~$f \colon A \to B$
و$g \colon B \to C$ تماثلين، كان~$(g \circ f) \colon A \to C$
تماثلًا كذلك.
\end{lem}
\begin{prob}
أقم البرهان على~\olref[sth][ordinals][iso]{isoscompose}.
\end{prob}
\begin{proof}
يُترك إثباتها للقارئ تمرينًا.
\end{proof}
\begin{cor}\ollabel{ordisoisequiv}
العلاقة~$\ordeq{X}{Y}$ علاقة تكافؤ.
\end{cor}
\begin{prop}\ollabel{ordisounique}
إذا تماثل الترتيبان الجيدان~$\tuple{A, <}$ و$\tuple{B, \lessdot}$،
لم يكن بينهما إلا تماثل واحد.
\end{prop}

\begin{proof}
لنفرض أن~$f$ و$g$ تماثلان~$A \to B$. وسبيلنا إلى إثبات تساويهما
الاستقراء الوارد في~\olref[wo]{propwoinduction}.
  نثبّت~$a\in A$، ونأخذ فرض الاستقراء
  $(\forall b < a)f(b) = g(b)$، ثم نثبّت~$x \in B$.

إن كان~$x \lessdot f(a)$، لزم~$f^{-1}(x) < a$، ثم
$g(f^{-1}(x)) \lessdot g(a)$، إذ إن~$f$ و$g$ تماثلان. ولما كان
$f^{-1}(x) < a$ أعطى فرض الاستقراء
$x =f(f^{-1}(x)) = g(f^{-1}(x))$؛ فيكون~$x \lessdot g(a)$.
وبالحجة نفسها، إن كان~$x \lessdot g(a)$ كان~$x \lessdot f(a)$.

إذن~$(\forall x \in B)(x \lessdot f(a) \liff x \lessdot g(a))$.
ومن~\olref[sfr][rel][ord]{prop:extensionality-strictlinearorders}
يلزم~$f(a) = g(a)$. وعليه
$(\forall a \in A)f(a) = g(a)$، بحسب~\olref[wo]{propwoinduction}.
\end{proof}\noindent
هذا أول ما يكشف عن ثبات بنية الترتيبات الجيدة. ولكي نستوفي البيان
نحتاج إلى ترميز آخر.

\begin{defn}
متى كان~$\tuple{A, <}$ ترتيبًا جيدًا و$a \in A$، وضعنا
$A_a = \Setabs{x \in A}{x < a}$. وتسمى~$A_a$ \emph{قطعة ابتدائية
حقيقية} من~$A$؛ أما~$A$ نفسها فنعدها قطعة ابتدائية غير حقيقية من~$A$.
ونرمز بـ$<_a$ إلى تقييد~$<$ بهذه القطعة، أي إلى
$\funrestrictionto{\mathord{<}}{A_a^2}$.
\end{defn}
\noindent
وبهذا الترميز نصوغ امتناع تماثل أي ترتيب جيد مع قطعة ابتدائية حقيقية
منه، ثم نبرهنه.

\begin{lem}\ollabel{wellordnotinitial}
إذا كان~$\tuple{A, <}$ ترتيبًا جيدًا و$a \in A$، لزم
$\ordneq{\tuple{A, <}}{\tuple{A_a, <_a}}$.
\end{lem}

\begin{proof}
لنفرض، طلبًا للخلف، أن~$f \colon A \to A_a$ تماثل. ولأن~$f$ تقابل
و$A_a \subsetneq A$، نأخذ~$b \in A$، بمقتضى
\olref[wo]{wo:strictorder}، أصغر !!{element} في~$A$ بحسب~$<$
مما يحقق~$b \neq f(b)$. وسنثبت
$(\forall x \in A)(x<b \liff x < f(b))$؛ إذ ينتج من ذلك، بمقتضى
\olref[sfr][rel][ord]{prop:extensionality-strictlinearorders}،
أن~$b = f(b)$، وهو الخلف المطلوب.

أما إذا كان~$x < b$، فإن~$x = f(x)$ بحكم اختيار~$b$. ثم
$f(x) < f(b)$ لحفظ~$f$ الترتيب، فيلزم~$x < f(b)$.

وأما إذا كان~$x < f(b)$، فإن~$f^{-1}(x) < b$ لأن~$f$ تماثل؛ فيكون
$f^{-1}(x) = x$ بحكم اختيار~$b$، ويلزم~$x < b$.
\end{proof}

وفي النتيجة التالية معنى يمكن تقريبه بالقول: إن «القطعة الابتدائية»
من التماثل تبقى تماثلًا.

\begin{lem}\ollabel{wellordinitialsegment}
ليكن~$\tuple{A, <}$ و$\tuple{B, \lessdot}$ ترتيبين جيدين. متى كان
$f \colon A \to B$ تماثلًا و$a \in A$، كان
$\funrestrictionto{f}{A_{a}} : A_a \to B_{f(a)}$ تماثلًا.
\end{lem}

\begin{proof}
من كون~$f$ تماثلًا نحصل على:
%Since $f$ is an isomorphism, $b < a$ iff $f(b) \lessdot f(a)$, so that
\begin{align*}
\funimage{f}{A_a} &= \funimage{f}{\Setabs{x \in A}{x < a}}\\
&= \Setabs{y \in B}{f^{-1}(y)<a} \\
&= \Setabs{y \in B}{y \lessdot f(a)} \\
&=B_{f(a)}
\end{align*}
وأما حفظ~$\funrestrictionto{f}{A_a}$ للترتيب فموروث من حفظ~$f$ له.
\end{proof}

ونثبت الآن، في نتيجتين، أن المقارنة بين الترتيبات الجيدة ممكنة دائمًا.

\begin{lem}\ollabel{lemordsegments}
ليكن~$\tuple{A, <}$ و$\tuple{B, \lessdot}$ ترتيبين جيدين. إذا كان
$\ordeq{\tuple{A_{a_1}, <_{a_1}}}{\tuple{B_{b_1}, \lessdot_{b_1}}}$
و$\ordeq{\tuple{A_{{a_2}}, <_{a_2}}}{\tuple{B_{{b_2}}, \lessdot_{b_2}}}$،
فإن~${a_1} < {a_2} \text{ إذا وفقط إذا كان } {b_1} \lessdot {b_2}$.
\end{lem}

\begin{proof}
نبرهن الاتجاه \emph{من اليسار إلى اليمين}؛ ويجري الآخر على مثاله.
لنفرض
$\ordeq{\tuple{A_{a_1}, <_{a_1}}}{\tuple{B_{b_1}, \lessdot_{b_1}}}$
و$\ordeq{\tuple{A_{{a_2}}, <_{a_2}}}{\tuple{B_{{b_2}}, \lessdot_{b_2}}}$،
وليكن~$f \colon A_{{a_2}} \to B_{{b_2}}$ التماثل الثاني. فإذا كان
${a_1} < {a_2}$، كان
$\ordeq{\tuple{A_{a_1}, <_{a_1}}}{\tuple{B_{f({a_1})}, \lessdot_{f({a_1})}}}$
بمقتضى~\olref{wellordinitialsegment}. ومن ثم
$\ordeq{\tuple{B_{b_1}, \lessdot_{b_1}}}{\tuple{B_{f({a_1})}, \lessdot_{f({a_1})}}}$،
فينتج~${b_1} = f({a_1})$ من~\olref{wellordnotinitial}. وأخيرًا
${b_1} \lessdot {b_2}$، إذ مدى~$f$ هو~$B_{{b_2}}$.
\end{proof}

\begin{thm}\ollabel{thm:woalwayscomparable}
كل ترتيبين جيدين يتماثل أحدهما مع قطعة ابتدائية من الآخر، ولا يلزم
أن تكون تلك القطعة حقيقية.
\end{thm}

\begin{proof}
نأخذ ترتيبين جيدين~$\tuple{A, <}$ و$\tuple{B, \lessdot}$، ونضع
بمسلّمة الفصل:
\[
f = \Setabs{\tuple{a, b} \in A \times B}{
\ordeq{\tuple{A_a, <_a}}{\tuple{B_b, \lessdot_b}}}.
\]
تقتضي~\olref{lemordsegments} أن~$a_1 < a_2$ إذا وفقط إذا كان
$b_1 \lessdot b_2$، لكل~$\tuple{a_1, b_1}, \tuple{a_2, b_2} \in f$.
فيكون~$f \colon \dom{f} \to \ran{f}$ تماثلًا.

وإذا كان~$a_2 \in \dom{f}$ و$a_1 < a_2$، لزم~$a_1 \in \dom{f}$
بمقتضى~\olref{wellordinitialsegment}؛ ومن ثم فإن~$\dom{f}$ قطعة
ابتدائية من~$A$. وكذلك~$\ran{f}$ قطعة ابتدائية من~$B$. ولنفرض
خلاف المطلوب، أي أن القطعتين \emph{حقيقيتان}. نأخذ~$a$ أصغر
!!{element} بحسب~$<$ في~$A \setminus \dom{f}$، فيكون~$\dom{f} = A_a$؛
ونأخذ~$b$ أصغر !!{element} بحسب~$\lessdot$ في~$B \setminus \ran{f}$،
فيكون~$\ran{f} = B_b$. إذن~$f \colon A_a \to B_b$ تماثل، فيلزم
$\tuple{a, b} \in f$، وهذا تناقض.
\end{proof}

\end{document}
